\label{chap.8}
\section{Summary}
In this project we looked at indirect prompt injection that spreads across the hops of a multi-agent LLM pipeline, including attacks split into pieces that look harmless one message at a time. We built MAPIS, a monitor that sits at a tap independent of the framework and remembers the session. It records where each item first appeared, gives a summary of that record and four signals to a fine-tuned DeBERTa-v3-small classifier, combines the result with rule signals and checks on action arguments, and responds with one of four tiers, a forensic trace and an operator dashboard.

We evaluated it on 2,304 held-out sessions from six benchmark sets, with five training seeds and five baselines. MAPIS reached $98.0\pm0.6$\% recall with $0.9\pm0.8$\% false positives and $98.5$\% accuracy at 63 ms per event. On multi-hop attacks it caught 100\% with 3.0\% false positives (at most 1.8\% in every seed with the calibrated cut-off), while a stateless model of the same size, limited to 5\% false positives, caught 4.3\%. Against an attacker trying to evade it, MAPIS caught 100\% at 0.4\% false positives, and with a live LLM agent it brought the attack success rate down from 15.3\% to 0\%. All four targets were met. The harder test sets showed that attacks which only change a value were a weak point; the tampering checks, added without retraining, raised recall on our team-written set from 40\% to 95\% and on attacks against the session record from 55.7\% to 80.4\%, with no new false positives.

\section{Answer to the Research Question}
What separates an attack from the same action requested by the user is the session memory. Without it, a detector either misses the fragment attacks or blocks the user's own requests. With it, a small and fast classifier can do both jobs.

\section{Limitations}
\begin{itemize}
\item The tampering checks were designed after reading the misses on the team-written set, so their 95\% on that set is not an independent estimate; a new hand-written set is needed to confirm it.
\item Attacks aimed at the session record reach 80.4\% recall but at 10.2\% false positives, above the 5\% target (mostly twins whose item was stored in memory); items the agent must decode (reversed or base64) are caught only 25\% of the time.
\item MultiHop and Adaptive are template-generated, with held-out families, wordings and tools and real carrier text.
\item At the default cut-off two of five seeds exceed 5\% false positives on multi-hop sessions; the validation-calibrated cut-off removes this in all five seeds.
\item MAPIS was trained on AgentDojo-style injections, and the live run uses them. The live agent, the LLM judge and NeMo used a small 7B local model.
\item On single-hop BIPIA there is no cross-hop evidence, so the stateless model is within seed variance of MAPIS in recall.
\item The four original rule signals add nothing measurable once the classifier has the session signals; only the tampering checks add recall.
\end{itemize}

\section{Scope for Further Research}
\begin{itemize}
\item Write a fresh hand-written test set to confirm the tampering checks, and extend them to items the agent must decode.
\item Tap agent-internal reads so that items introduced by an agent can be traced, and evaluate on real agent traffic.
\item Run larger live agents and stronger LLM judges as baselines.
\item Extend the tap to multimodal content and to additional frameworks, and harden the system for production use.
\end{itemize}
